\documentclass[11pt]{article}

\usepackage[T1]{fontenc}
\usepackage[utf8]{inputenc}
\usepackage{lmodern}
\usepackage{microtype}
\usepackage[margin=1in]{geometry}
\usepackage{amsmath,amssymb,amsthm,mathtools}
\usepackage{booktabs,array,longtable}
\usepackage{enumitem}
\usepackage{xcolor}
\usepackage{hyperref}
\usepackage[nameinlink,noabbrev]{cleveref}
\usepackage{listings}
\usepackage{fancyhdr}

\hypersetup{
  colorlinks=true,
  linkcolor=blue!55!black,
  citecolor=blue!55!black,
  urlcolor=blue!55!black,
  pdftitle={Pappus's centroid theorem in intuitionistic primitive recursive arithmetic},
  pdfauthor={OpenAI}
}

\pagestyle{fancy}
\setlength{\headheight}{14pt}
\fancyhf{}
\lhead{Pappus's centroid theorem in iPRA}
\rhead{Constructive proof package}
\cfoot{\thepage}

\newtheorem{theorem}{Theorem}[section]
\newtheorem{lemma}[theorem]{Lemma}
\newtheorem{proposition}[theorem]{Proposition}
\newtheorem{corollary}[theorem]{Corollary}
\theoremstyle{definition}
\newtheorem{definition}[theorem]{Definition}
\newtheorem{construction}[theorem]{Construction}
\newtheorem{remark}[theorem]{Remark}
\newtheorem{example}[theorem]{Example}

\newcommand{\N}{\mathbb N}
\newcommand{\Z}{\mathbb Z}
\newcommand{\Q}{\mathbb Q}
\newcommand{\R}{\mathbb R}
\newcommand{\D}{\mathbb D}
\newcommand{\ipra}{\mathsf{iPRA}}
\newcommand{\PR}{\mathsf{PR}}
\newcommand{\Eq}{\mathrel{=_{\!\mathsf R}}}
\newcommand{\Le}{\mathrel{\le_{\!\mathsf R}}}
\newcommand{\dist}{\operatorname{dist}}
\newcommand{\Area}{\operatorname{Area}}
\newcommand{\Vol}{\operatorname{Vol}}
\newcommand{\Len}{\operatorname{Len}}
\newcommand{\Mom}{\operatorname{Mom}}
\newcommand{\Surf}{\operatorname{Surf}}
\newcommand{\code}[1]{\texttt{\detokenize{#1}}}
\newcommand{\abs}[1]{\left|#1\right|}
\newcommand{\norm}[1]{\left\lVert#1\right\rVert}
\newcommand{\set}[1]{\left\{#1\right\}}
\newcommand{\brak}[1]{\left[#1\right]}
\newcommand{\ceil}[1]{\left\lceil #1\right\rceil}

\lstdefinelanguage{Lean}{
  morekeywords={def,theorem,lemma,namespace,end,by,have,show,from,import,open,section,variable,where,if,then,else,let,in,match,with,case,induction,exact,refine,simp,rw,constructor,structure,noncomputable,axiom},
  sensitive=true,
  morecomment=[l]{--},
  morecomment=[s]{/-}{-/},
  morestring=[b]"
}
\lstset{
  language=Lean,
  basicstyle=\ttfamily\small,
  columns=fullflexible,
  breaklines=true,
  frame=single,
  showstringspaces=false,
  backgroundcolor=\color{black!2}
}

\title{\textbf{Pappus's Centroid Theorem in Intuitionistic Primitive Recursive Arithmetic}\\
       \large Exact polygonal and grid-cell identities, explicit constructive limits, and a Lean formalization route}
\author{Research report prepared for Vladimir Reshetnikov}
\date{15 November 2026}

\begin{document}
\maketitle

\begin{abstract}
Pappus's two centroid theorems assert that a surface of revolution has area
``generating arc length times the distance traveled by the arc centroid,'' and
that a solid of revolution has volume ``generating area times the distance
traveled by the lamina centroid.''  An unrestricted statement about arbitrary
real curves, arbitrary measurable laminas, and completed real integrals is not
a first-order sentence of intuitionistic primitive recursive arithmetic
(\(\ipra\)).  We therefore give an exact represented-data formulation and a
fully constructive proof.

The proof has a finite core.  For one meridional line segment, an explicit squeeze by inscribed and
circumscribed polyhedral frusta gives
\[
  S=\tau\,\frac{r_0+r_1}{2}\,\ell,
\]
where \(\tau=2\pi\), \(r_0,r_1\) are the endpoint radii, and \(\ell\) is the
segment length.  The factor multiplying \(\tau\) is exactly the radial first
moment of the segment.  For one rectangular grid cell
\([a,b]\times[c,d]\), rotation gives an annular cylinder and
\[
  V=\frac{\tau}{2}(b^2-a^2)(d-c)
   =\tau\,\frac{a+b}{2}(b-a)(d-c),
\]
the latter factor being the cell's radial first moment.  Finite additivity
proves both theorems for polygonal chains and finite grid laminas.

General constructive curves and laminas are presented by primitive-recursive
polygonal approximants or nested inner/outer dyadic cell unions, together with
explicit Cauchy or boundary-cover moduli.  The exact finite identities pass to
the represented-real limits with an explicit precision shift.  Division by
arc length or area is used only with a supplied positive lower witness.  The
circle constant is itself given by a primitive-recursive regular-polygon
presentation; no least-upper-bound axiom, measure-theoretic change of
variables, countable choice, or excluded middle for arbitrary propositions is
used.  Every induction reduces to quantifier-free induction on coded finite
objects.  We conclude with a module-by-module Lean plan tied to the current
\code{ConciseIPRA} development, a no-\code{sorry} Lean finite-algebra scaffold (not locally kernel-checked), and
an executable exact-rational checker.
\end{abstract}

\tableofcontents

\section{The exact constructive claim}
\label{sec:claim}

The expository statement normally called Pappus's centroid theorem has two
parts.  In coordinates where the rotation axis is the \(z\)-axis and the
meridional half-plane is \(r\geq0\), write \(\tau=2\pi\).

\begin{enumerate}[label=\textbf{(S\arabic*)},leftmargin=2.7em]
\item If a curve of length \(L\) has radial centroid coordinate \(\bar r\),
then the area of the swept surface is
\[
  S=\tau L\bar r.
\]
The centroid travels a circle of length \(d=\tau\bar r\), so \(S=Ld\).

\item If a lamina of area \(A\) has radial centroid coordinate \(\bar r\),
then the volume of the swept solid is
\[
  V=\tau A\bar r.
\]
Again \(d=\tau\bar r\), so \(V=Ad\).
\end{enumerate}

The usual hypotheses are that the generator lies on one side of the axis and
that the rotation does not cover a two-dimensional or three-dimensional part
more than once.  Contact with the axis on a zero-content boundary set is
harmless; this includes the usual cone, cylinder, and sphere examples.

The phrase ``prove this in \(\ipra\)'' needs a precise first-order reading.
The theory has natural numbers and a symbol for every fixed primitive-recursive
function; it does not quantify over arbitrary functions \(\R\to\R^2\),
arbitrary subsets of the plane, or an impredicative completed real field.
The correct formulation is therefore a theorem schema for each fixed
primitive-recursive presentation.

\begin{theorem}[Constructive Pappus theorem: surface version]
\label{thm:surface-main}
Let \(\mathcal C\) be a fixed primitive-recursive polygonal presentation of a
rectifiable meridional curve.  At stage \(n\), it produces a finite rational
polygonal chain \(P_n\) in \(r\geq0\).  Suppose the presentation supplies:
\begin{enumerate}[label=(\roman*),leftmargin=2.3em]
\item regular real names for the polygonal length \(L_n\) and radial first
moment \(M_n\);
\item a primitive-recursive Cauchy modulus for \((L_n)\) and \((M_n)\);
\item a positive rational lower witness for the limiting length \(L\);
\item a finite certificate excluding nontrivial retracing, or else the
interpretation of area with multiplicity.
\end{enumerate}
Let \(L=\lim_nL_n\), \(M=\lim_nM_n\), let
\(\bar r=M/L\), and define the swept surface area \(S\) as the limit of the
areas of the revolved polygonal chains.  Then \(\ipra\) proves equality of the
corresponding regular real names:
\[
  S\Eq \tau M\Eq \tau L\bar r.
\]
Consequently, with \(d=\tau\bar r\), \(S\Eq Ld\).
\end{theorem}

\begin{theorem}[Constructive Pappus theorem: volume version]
\label{thm:volume-main}
Let \(\mathcal D\) be a fixed primitive-recursive Jordan-lamina presentation
in a bounded rectangle \(0\leq r\leq R\).  At precision stage \(n\), it
produces finite unions \(I_n\subseteq O_n\) of pairwise interior-disjoint
dyadic grid cells, interpreted as inner and outer approximants, such that
\[
 I_n\subseteq I_{n+1}\subseteq O_{n+1}\subseteq O_n,
 \qquad
 \Area(O_n\setminus I_n)\leq2^{-n}.
\]
Suppose also that the limiting area has a supplied positive rational lower
witness.  Let \(A\), \(M\), and \(V\) be the common limits of area, radial
first moment, and swept volume; define \(\bar r=M/A\).  Then \(\ipra\) proves
\[
  V\Eq \tau M\Eq \tau A\bar r.
\]
Consequently, with \(d=\tau\bar r\), \(V\Eq Ad\).
\end{theorem}

\begin{corollary}[Classical reading]
\label{cor:classical-reading}
In ordinary constructive analysis, a simple piecewise \(C^1\) curve with an
explicit modulus of differentiability produces a presentation satisfying
\cref{thm:surface-main}.  A bounded Jordan lamina with an explicit
primitive-recursive boundary-cover modulus produces a presentation satisfying
\cref{thm:volume-main}.  Under classical logic, the represented quantities
coincide with the usual arc length, line centroid, area centroid, surface area,
and volume.  Hence the two theorems recover the customary Pappus formulas for
all such effective data.
\end{corollary}

\begin{remark}[What is and is not being restricted]
The mathematical identity is not weakened to rational polygons only.  The
finite polygon and grid identities are the exact kernel; explicit limits then
cover regular curves and regions.  The genuine restriction is proof-theoretic:
a curve or lamina must be supplied by algorithms and moduli.  An assertion
about an arbitrary, unnamed, noncomputable set of reals has no direct
first-order meaning in \(\ipra\).
\end{remark}

\section{Arithmetic coding and regular real names}
\label{sec:reals}

\subsection{Rational and finite-object codes}

A signed rational is coded by a sign bit and a pair \((a,b)\) with
\(a\in\N\), \(b>0\).  Equality and order are decided by cross multiplication.
Dyadics use denominators \(2^e\).  Finite lists are coded by a standard pairing
function together with primitive-recursive length and projection functions.
The following operations are primitive recursive:
\begin{itemize}[leftmargin=2em]
\item rational addition, multiplication, negation, comparison, and reciprocal
with a supplied positive-denominator witness;
\item bounded sums and products, finite maxima and minima, list maps and folds;
\item integer square, absolute difference, bounded division and remainder;
\item decoding polygon vertices and finite sets of dyadic grid cells.
\end{itemize}
Every rational identity below expands to an integer identity.  Every rational
inequality expands to an integer inequality after multiplying by explicitly
positive denominators.

\subsection{Regular interval names}

\begin{definition}[Regular real name]
A regular real name \(X\) is a pair of primitive-recursive rational functions
\(L_X,U_X:\N\to\Q\) satisfying, for every \(k\),
\begin{align*}
 L_X(k)&\leq L_X(k+1)\leq U_X(k+1)\leq U_X(k),\\
 U_X(k)-L_X(k)&\leq2^{-k}.
\end{align*}
Write \(c_X(k)=(L_X(k)+U_X(k))/2\).
\end{definition}

We use the arithmetic equality convention
\begin{equation}
 X\Eq Y
 \quad:\Longleftrightarrow\quad
 \forall k\;\abs{c_X(k+3)-c_Y(k+3)}\leq2^{-k}.
 \label{eq:real-eq}
\end{equation}
The three-place precision shift makes reflexivity, symmetry, transitivity, and
compatibility with interval arithmetic immediate.  Order is defined similarly:
\begin{equation}
 X\Le Y
 \quad:\Longleftrightarrow\quad
 \forall k\;c_X(k+3)\leq c_Y(k+3)+2^{-k}.
 \label{eq:real-le}
\end{equation}
At each fixed precision these are decidable rational inequalities.

\begin{definition}[Positive lower witness]
A name \(X\) is strongly positive when its data include rational
\(q>0\) and an index \(p\) such that \(q\leq L_X(p)\).  Nesting makes \(q\) a
permanent lower bound.  Reciprocals and cancellation are used only with such a
witness.
\end{definition}

Addition and multiplication are performed by interval arithmetic and a
precision shift large enough to absorb known magnitude bounds.  A compatible
but nonnested family of intervals is normalized by
\[
 \widehat L(k)=\max_{i\leq k}L(i),
 \qquad
 \widehat U(k)=\min_{i\leq k}U(i).
\]
All these transformations are primitive recursive.

\subsection{Square roots by bounded bisection}

For a nonnegative rational \(q\), choose an integer bound \(B\geq1+q\).
Starting with \([0,B]\), perform a predetermined number of bisection steps.
At a midpoint \(m\), the comparison \(m^2\leq q\) is decidable, so one retains
\([m,u]\) or \([\ell,m]\).  After
\[
 N(k,B)=k+2+\min\set{j:B\leq2^j}
\]
steps the width is at most \(2^{-k-2}\).  The least \(j\leq B+1\) with the
displayed property is obtained by bounded search.  Thus \(\sqrt q\) has a
primitive-recursive regular interval name, and \(\ipra\) proves the invariant
\[
 0\leq\ell_N\leq\sqrt q\leq u_N,
 \qquad
 u_N-\ell_N\leq2^{-k-2},
\]
using quantifier-free induction on the fixed bisection counter.  Segment
lengths of rational polygonal chains therefore have uniform regular names.

\subsection{Explicit limits}

\begin{lemma}[Primitive-recursive Cauchy completion]
\label{lem:pr-limit}
Suppose \(X_n\) is a primitive-recursive family of regular names and a
primitive-recursive function \(\mu\) satisfies
\[
 m,n\geq\mu(k)\quad\Longrightarrow\quad
 \abs{X_m-X_n}\leq2^{-k}
\]
in the represented-real sense.  Then the diagonal construction
\[
 c_X(k)=c_{X_{\mu(k+4)}}(k+4)
\]
with interval radius \(2^{-k-1}\), followed by running intersection, is a
regular name for a limit \(X\).  If \(Y_n\Eq X_n\) at every stage and
\(Y_n\to Y\) with an explicit modulus, then \(X\Eq Y\).
\end{lemma}

\begin{proof}
Given precision \(k\), every two selected diagonal stages lie beyond
\(\mu(k+4)\), and the internal approximation error of each selected name is at
most \(2^{-k-4}\).  Three applications of the triangle inequality give a
total error below \(2^{-k-1}\).  All indices are explicit primitive-recursive
terms.  The uniqueness assertion uses the same calculation with the exact
stagewise equality inserted in the middle.  At fixed \(k\), the calculation is
a finite rational inequality; the only iteration is over the coded finite
sum of error terms.
\end{proof}

\begin{lemma}[Compatibility with a bounded factor]
\label{lem:bounded-factor}
If \(X_n\to X\) with modulus \(\mu\), and \(C\) is a regular name with a
known bound \(\abs C\leq B\), then \(CX_n\to CX\) with modulus
\(k\mapsto\mu(k+\ceil{\log_2(B+1)}+3)\).  The logarithm notation denotes a
bounded integer search, not a real logarithm.
\end{lemma}

\begin{proof}
Choose the stage so that \(\abs{X_n-X}\leq2^{-k}/(B+1)\).  Interval-product
rounding consumes two additional bits.  Multiplying the rational inequalities
and clearing positive denominators proves the claim.
\end{proof}

\section{A primitive-recursive circle constant package}
\label{sec:circle}

The Pappus identities use one geometric constant: the length of a full turn
at unit radius.  We denote it by \(\tau\); later \(\pi=\tau/2\).  It is useful
to construct \(\tau\) directly rather than assume a completed transcendental
real.

\subsection{Regular polygon brackets}

Let \(N_m=2^{m+2}\).  Define regular real names
\begin{align*}
 c_0&=\sqrt{\tfrac12},
 &c_{m+1}&=\sqrt{\frac{1+c_m}{2}},\\
 s_m&=\sqrt{1-c_m^2},
 &p_m&=2N_ms_m,
 &q_m&=\frac{p_m}{c_m}.
\end{align*}
Geometrically, \(c_m\) is the apothem of an inscribed regular \(N_m\)-gon in
the unit circle, \(p_m\) is its perimeter, and \(q_m\) is the perimeter of
the corresponding circumscribed polygon.  The reciprocal is constructive
because \(c_m\geq c_0>1/2\).

\begin{lemma}[Elementary polygon estimates]
\label{lem:polygon-estimates}
For every \(m\),
\begin{align}
 \tfrac12&\leq c_m\leq1,
 &1-c_m&\leq2^{-m},\label{eq:c-bound}\\
 p_m&\leq p_{m+1}\leq q_{m+1}\leq q_m\leq8,
 &q_m-p_m&\leq8\cdot2^{-m}.
 \label{eq:perim-gap}
\end{align}
If
\[
 a_m=\frac12p_mc_m,
 \qquad
 b_m=\frac12q_m,
\]
are the inscribed and circumscribed polygon areas, then
\begin{equation}
 a_m\leq a_{m+1}\leq b_{m+1}\leq b_m,
 \qquad
 b_m-a_m\leq8\cdot2^{-m}.
 \label{eq:area-gap}
\end{equation}
Moreover the common perimeter limit is twice the common area limit.
\end{lemma}

\begin{proof}
From \(c_{m+1}^2=(1+c_m)/2\),
\[
 1-c_{m+1}
 =\frac{1-c_m}{2(1+c_{m+1})}
 \leq\frac{1-c_m}{2},
\]
which proves \eqref{eq:c-bound} by induction.  Half-angle similarity gives
\(p_m=q_mc_m\).  The doubled polygon satisfies
\[
 \frac{p_{m+1}}{p_m}=\frac1{c_{m+1}}\geq1,
 \qquad
 \frac{q_{m+1}}{q_m}=\frac{2c_m}{1+c_m}\leq1.
\]
The circumscribed square has perimeter \(q_0=8\), hence \(q_m\leq8\), and
\[
 q_m-p_m=q_m(1-c_m)\leq8\cdot2^{-m}.
\]
The area formula is the finite sum of triangle areas, so
\(a_m=p_mc_m/2\) and \(b_m=q_m/2\).  In fact
\[
 a_{m+1}=\frac{p_m}{2}\geq\frac{p_mc_m}{2}=a_m,
 \qquad b_{m+1}\leq b_m,
 \qquad a_{m+1}\leq b_{m+1}.
\]
Therefore
\[
 b_m-a_m=\frac{q_m(1-c_m^2)}2
 \leq q_m(1-c_m)
 \leq8\cdot2^{-m}.
\]
Finally
\[
 \abs{p_m-2a_m}=p_m(1-c_m)\leq8\cdot2^{-m},
 \qquad 2b_m=q_m.
\]
Both perimeter brackets and area brackets have vanishing explicit gaps, so
\cref{lem:pr-limit} gives common limits \(\tau\) and \(\pi\), and the last two
relations give \(\tau=2\pi\).  Every displayed comparison is algebraic in
square-root names and is proved at fixed precision by denominator clearing.
\end{proof}

\begin{definition}[Circle package]
Let \(\tau\) be the regular real defined by the nested perimeter brackets
\([p_m,q_m]\), after the precision shift \(m=k+4\).  Define
\(\pi=\tau/2\).  A circle of radius \(r\geq0\) has circumference \(\tau r\),
and a disk of radius \(r\) has area \(\tau r^2/2\), by scaling the polygon
approximants.
\end{definition}

The package proves \(5<\tau<8\), more than enough for every later error budget.
No trigonometric function is needed in the object theory.


\section{The finite surface identity}
\label{sec:finite-surface}

This section proves the geometric kernel of the first theorem.  Everything is
reduced to finite Euclidean algebra and the regular-polygon brackets of
\cref{sec:circle}.  In particular, the proof does not assume a pre-existing
surface integral.

\subsection{The first moment of one segment}

Let
\[
 u_0=(r_0,z_0),\qquad u_1=(r_1,z_1),\qquad r_0,r_1\geq0,
\]
and put
\[
 \Delta r=r_1-r_0,\qquad \Delta z=z_1-z_0,
 \qquad \ell=\sqrt{(\Delta r)^2+(\Delta z)^2}.
\]
The segment is parametrized affinely by
\[
 u(t)=((1-t)r_0+tr_1,(1-t)z_0+tz_1),\qquad0\leq t\leq1.
\]

\begin{lemma}[Finite midpoint moment]
\label{lem:segment-moment}
For every positive integer $N$, divide the segment into $N$ equal pieces and
place the mass of each piece at its midpoint.  The radial first moment of this
finite mass distribution is exactly
\[
 M_N=\ell\,\frac{r_0+r_1}{2}.
\]
It is independent of $N$.  Hence the constructive line moment and radial
centroid of the nondegenerate segment are
\[
 M_\sigma=\ell\,\frac{r_0+r_1}{2},
 \qquad
 \bar r_\sigma=\frac{r_0+r_1}{2}.
\]
\end{lemma}

\begin{proof}
The $j$th midpoint, for $0\leq j<N$, has parameter
$t_j=(2j+1)/(2N)$ and radial coordinate
$r_j=r_0+t_j(r_1-r_0)$.  Each piece has length $\ell/N$.  Thus
\begin{align*}
 M_N
 &=\frac{\ell}{N}\sum_{j=0}^{N-1}
   \left(r_0+\frac{2j+1}{2N}(r_1-r_0)\right)\\
 &=\ell r_0+\frac{\ell(r_1-r_0)}{2N^2}
   \sum_{j=0}^{N-1}(2j+1).
\end{align*}
The finite odd-number identity
$\sum_{j<N}(2j+1)=N^2$ gives the result.  That identity is proved by
quantifier-free induction.  No limiting argument is needed for the affine
function $r(t)$.
\end{proof}

\subsection{Polyhedral brackets for the revolved segment}

Rotate the segment about the $z$-axis.  The familiar picture is a conical
frustum, with the cylinder and cone as limiting cases.  To define its area
constructively, use the same regular $N_m$-gons as in the circle package.
Write $c_m$ and $s_m$ for their half-angle cosine and sine, and $p_m,q_m$ for
the unit-circle inscribed and circumscribed perimeters.

For the inner approximation, put an inscribed regular $N_m$-gon in each
endpoint circle and join corresponding vertices.  One face has parallel base
lengths $2r_0s_m$ and $2r_1s_m$.  In coordinates centered on the midpoint of
one polygon edge, the perpendicular displacement between the two base lines
is
\[
 h_m=\sqrt{(\Delta z)^2+c_m^2(\Delta r)^2}.
\]
The face is an ordinary planar trapezoid.  Summing its area over the $N_m$
faces gives
\begin{equation}
 S_m^-=\frac{p_m}{2}(r_0+r_1)h_m.
 \label{eq:inner-frustum}
\end{equation}

For the outer approximation, put a circumscribed regular $N_m$-gon around each
endpoint circle and again join corresponding vertices.  A tangent side at
radius $r$ has length $2r\,s_m/c_m$.  The perpendicular displacement between
corresponding tangent sides is exactly $\ell$.  Hence
\begin{equation}
 S_m^+=\frac{q_m}{2}(r_0+r_1)\ell.
 \label{eq:outer-frustum}
\end{equation}
All coordinates involved here are obtained from rational input by the square
root and field operations already represented in \cref{sec:reals}.

\begin{lemma}[Effective frustum squeeze]
\label{lem:frustum-squeeze}
The inner areas increase, the outer areas decrease, and
\[
 S_m^-\leq S_{m+1}^-\leq S_{m+1}^+\leq S_m^+.
\]
Moreover
\begin{equation}
 0\leq S_m^+-S_m^-
 \leq4(r_0+r_1)(\ell+\abs{\Delta r})2^{-m}.
 \label{eq:frustum-gap}
\end{equation}
The two sequences therefore define one regular real, and that real is
\begin{equation}
 \Surf(\sigma)=\tau\,\frac{r_0+r_1}{2}\,\ell.
 \label{eq:segment-surface}
\end{equation}
\end{lemma}

\begin{proof}
The sequences $p_m,c_m$ increase and $q_m$ decreases.  Since
$h_m^2=(\Delta z)^2+c_m^2(\Delta r)^2$, the sequence $h_m$ also increases and
never exceeds $\ell$.  This proves the nesting.

The reverse triangle inequality, applied to the two vectors
$(\Delta z,c_m\Delta r)$ and $(\Delta z,\Delta r)$, gives the division-free
bound
\[
 0\leq\ell-h_m\leq(1-c_m)\abs{\Delta r}
 \leq2^{-m}\abs{\Delta r}.
\]
Using $p_m\leq q_m\leq8$ and $q_m-p_m\leq8\cdot2^{-m}$,
\begin{align*}
 S_m^+-S_m^-
 &=\frac{r_0+r_1}{2}\bigl(q_m\ell-p_mh_m\bigr)\\
 &=\frac{r_0+r_1}{2}
   \bigl((q_m-p_m)\ell+p_m(\ell-h_m)\bigr)\\
 &\leq4(r_0+r_1)(\ell+\abs{\Delta r})2^{-m}.
\end{align*}
A rational upper bound for the constant in front is computable from the input
codes, so a primitive-recursive precision shift turns these brackets into a
regular name.

Finally $p_m\to\tau$ and $h_m\to\ell$ with the displayed moduli.  Passing to
the common limit in \eqref{eq:inner-frustum} gives
\eqref{eq:segment-surface}.  More explicitly, for every requested $k$, choose
$m$ so that the right side of \eqref{eq:frustum-gap} is at most $2^{-(k+2)}$
and use either bracket; this is the object-theoretic proof of real equality.
\end{proof}

\begin{remark}[All degeneracies are covered]
No division by $\Delta r$ occurred.  Thus the same proof includes a cylinder
($\Delta r=0$), a cone ($r_0=0$ or $r_1=0$), a segment on the axis
($r_0=r_1=0$, giving zero swept area), and an endpoint touching the axis.
This is one reason the polygonal squeeze is preferable to the customary proof
that extends the sloping side to the cone apex and unrolls an annular sector.
The latter is valid intuition, but it introduces a case split and an implicit
sector-area theorem that would itself need constructive justification.
\end{remark}

Combining \cref{lem:segment-moment,lem:frustum-squeeze} immediately yields the
one-segment form of Pappus's theorem.

\begin{proposition}[Segment Pappus identity]
\label{prop:segment-pappus}
For every meridional segment in $r\geq0$,
\[
 \Surf(\sigma)=\tau M_\sigma.
\]
If $\ell>0$, this is
\[
 \Surf(\sigma)=\tau\ell\bar r_\sigma.
\]
\end{proposition}

\subsection{Finite polygonal chains}

Let $P=(u_0,\ldots,u_s)$ be a finite polygonal chain in $r\geq0$.  For its
$j$th segment write $\ell_j$ for its length and $r_j,r_{j+1}$ for its endpoint
radii.  Define
\begin{align}
 L(P)&=\sum_{j<s}\ell_j,\label{eq:poly-length}\\
 M(P)&=\sum_{j<s}\frac{r_j+r_{j+1}}2\ell_j,
 \label{eq:poly-moment}\\
 S(P)&=\sum_{j<s}\Surf([u_j,u_{j+1}]).
 \label{eq:poly-surface}
\end{align}
The last definition counts area with multiplicity.  It is always meaningful,
even for a self-intersecting chain.

\begin{theorem}[Finite polygonal Pappus theorem]
\label{thm:polygon-pappus}
For every finite polygonal chain $P$ in $r\geq0$,
\begin{equation}
 S(P)=\tau M(P).
 \label{eq:polygon-pappus}
\end{equation}
If $L(P)>0$ and $\bar r(P)=M(P)/L(P)$, then
\begin{equation}
 S(P)=\tau L(P)\bar r(P).
 \label{eq:polygon-centroid}
\end{equation}
If the chain is simple away from possible axis endpoints, then $S(P)$ is the
ordinary geometric area of its surface of revolution; otherwise it is the
ordinary area counted with multiplicity.
\end{theorem}

\begin{proof}
Sum \cref{prop:segment-pappus} over the finitely many segments and distribute
multiplication by $\tau$ over the finite sum.  This proves
\eqref{eq:polygon-pappus}.  When $L(P)>0$, the positive reciprocal construction
gives $L(P)\bar r(P)=M(P)$, and \eqref{eq:polygon-centroid} follows.

For the final assertion, a point of a revolved segment with positive radius is
determined by its meridional pair $(r,z)$ and its angular position.  Two
nonadjacent pieces can therefore overlap on a surface patch only if their
meridional segments overlap on a nontrivial subsegment.  Simplicity excludes
this.  Adjacent pieces share only an endpoint circle, and points on the axis
collapse only a one-dimensional angular fiber to a point; both have zero
surface content by explicit polygonal covers.  Finite additivity therefore
identifies the sum with geometric area.  No quotient by an undecidable
``almost everywhere'' relation is required.
\end{proof}

\section{The finite volume identity}
\label{sec:finite-volume}

The volume theorem has an even simpler finite kernel.  Let
\[
 C=[a,b]\times[c,d],\qquad0\leq a\leq b,\qquad c\leq d,
\]
be one meridional rectangle.  Its area and radial first moment are
\begin{equation}
 A(C)=(b-a)(d-c),
 \qquad
 M(C)=\frac{a+b}{2}(b-a)(d-c).
 \label{eq:cell-area-moment}
\end{equation}
The moment formula can again be proved by finite midpoint sums: on every
uniform radial subdivision the average of the cell-center radii is
$(a+b)/2$.

Rotating $C$ gives an annular cylinder.  By the circle package, its volume is
\begin{equation}
 V(C)=\frac{\tau}{2}(b^2-a^2)(d-c).
 \label{eq:cell-volume}
\end{equation}
Internally this is defined by subtracting the regular polygonal-prism names
for the two cylinders.  Thus \eqref{eq:cell-volume} is not an appeal to
Cavalieri's principle over arbitrary sets; it is a limit of finite prism
identities with the modulus inherited from \eqref{eq:area-gap}.

\begin{proposition}[Cell Pappus identity]
\label{prop:cell-pappus}
For every such cell,
\[
 V(C)=\tau M(C).
\]
\end{proposition}

\begin{proof}
The denominator-cleared calculation is finite:
\begin{align*}
 V(C)
 &=\frac{\tau}{2}(b-a)(b+a)(d-c)\\
 &=\tau\frac{a+b}{2}(b-a)(d-c)
 =\tau M(C).
\end{align*}
The only division by $2$ is multiplication by the fixed rational $1/2$.
\end{proof}

Let $G=(C_0,\ldots,C_{t-1})$ be a finite list of grid cells with pairwise
disjoint interiors.  Define
\[
 A(G)=\sum_{i<t}A(C_i),\qquad
 M(G)=\sum_{i<t}M(C_i),\qquad
 V(G)=\sum_{i<t}V(C_i).
\]

\begin{theorem}[Finite grid-lamina Pappus theorem]
\label{thm:grid-pappus}
For every finite grid lamina $G$,
\[
 V(G)=\tau M(G).
\]
If $A(G)>0$ and $\bar r(G)=M(G)/A(G)$, then
\[
 V(G)=\tau A(G)\bar r(G).
\]
The sum is the geometric volume of the rotated union, not merely a volume with
multiplicity.
\end{theorem}

\begin{proof}
Sum \cref{prop:cell-pappus}.  Interior-disjoint meridional cells rotate to
interior-disjoint annular cylinders: equality of a rotated point with positive
radius forces equality of its $(r,z)$ pair.  Shared faces and circular edges
have zero volume, witnessed by boxes of arbitrarily small total rational
volume.  Hence finite additivity applies.  The centroid formula follows from
$A(G)\bar r(G)=M(G)$ using the supplied positivity of $A(G)$.
\end{proof}

\section{Explicit limits and proof of the two main theorems}
\label{sec:limits}

We now pass from finite generators to the represented objects of
\cref{thm:surface-main,thm:volume-main}.  The passage is constructive because
all rates are part of the input.

\subsection{A uniform limit-transfer lemma}

The following elementary lemma is used repeatedly.  Its formulation is chosen
so that every witness is a primitive-recursive term.

\begin{lemma}[Transfer of a finite identity through explicit limits]
\label{lem:identity-limit}
Let $x_n,y_n$ be sequences of regular reals with fixed primitive-recursive
Cauchy moduli, and suppose $x_n=c\,y_n$ for every $n$, where $c$ is a fixed
regular real with a rational bound $\abs c\leq C$.  Let $x=\lim x_n$ and
$y=\lim y_n$ be the diagonal regular names of \cref{lem:pr-limit}.  Then
$x=c\,y$.

More quantitatively, if
\[
 \abs{x-x_n}\leq2^{-(k+2)},\qquad
 \abs{y-y_n}\leq2^{-(k+2+B)}
\]
and $C\leq2^B$, then $\abs{x-cy}\leq2^{-k}$.
\end{lemma}

\begin{proof}
For the chosen $n$,
\[
 \abs{x-cy}
 \leq\abs{x-x_n}+\abs{x_n-cy_n}+\abs c\,\abs{y_n-y}
 \leq2^{-(k+2)}+0+2^B2^{-(k+2+B)}
 \leq2^{-(k+1)}.
\]
The remaining slack pays for extracting rational endpoints from the regular
names.  All comparisons are decidable rational inequalities after a fixed
precision shift.  The stage $n$ is the maximum of the two supplied moduli, so
it is primitive recursive.
\end{proof}

\begin{lemma}[Positive division]
\label{lem:positive-division}
Let $X,Y$ be regular reals, and suppose the data include a positive rational
$\lambda$ and a stage after which $L_X(k)\geq\lambda$.  Then the reciprocal
$1/X$ and quotient $Y/X$ have primitive-recursive regular names.  They satisfy
$X(Y/X)=Y$.  No disjunction $X=0\vee X\neq0$ is used.
\end{lemma}

\begin{proof}
On the compact rational interval $[\lambda,U_X(0)]$, reciprocal has the
explicit modulus
\[
 \abs{u^{-1}-v^{-1}}=\frac{\abs{u-v}}{uv}
 \leq\lambda^{-2}\abs{u-v}.
\]
Choose a power of two exceeding $\lambda^{-2}$, refine the input name by the
corresponding precision shift, and reverse the interval endpoints.  The field
identity follows at every requested precision by cross multiplication.  The
construction is primitive recursive because $\lambda$ is supplied, rather
than found by an unbounded search.
\end{proof}

\subsection{Surface presentations}

For a polygonal presentation $n\mapsto P_n$, let
\[
 L_n=L(P_n),\qquad M_n=M(P_n),\qquad S_n=S(P_n).
\]
By \cref{thm:polygon-pappus},
\begin{equation}
 S_n=\tau M_n
 \qquad\text{for every }n.
 \label{eq:stage-surface}
\end{equation}
The presentation supplies Cauchy moduli for $L_n$ and $M_n$.  Since
$\abs\tau<8$, equation \eqref{eq:stage-surface} itself gives a Cauchy modulus
for $S_n$: a modulus for $M_n$ at precision $k+4$ is one for $S_n$ at
precision $k$.  Thus all three diagonal limits are regular names.

\begin{proof}[Proof of \cref{thm:surface-main}]
Apply \cref{lem:identity-limit} to \eqref{eq:stage-surface}, with $c=\tau$ and
$C=8$.  This gives
\[
 S=\tau M.
\]
The positive lower witness for $L$ permits the definition
$\bar r=M/L$ by \cref{lem:positive-division}; hence $L\bar r=M$.  Associativity
of represented-real multiplication gives
\[
 S=\tau M=\tau(L\bar r)=\tau L\bar r.
\]
With $d=\tau\bar r$, this is $S=Ld$.

At the object-theory level, the phrase ``apply the limit lemma'' expands as
follows.  Given a precision variable $k$, compute
\[
 n=\max\{\mu_S(k+3),\mu_M(k+7)\}.
\]
Open the finite certificate for $P_n$, invoke the segment identities, fold the
finite sum, and use the two Cauchy bounds and $\tau<8$.  After denominators are
cleared, the goal is a quantifier-free numerical inequality.  Its universal
closure in $k$ is therefore an ordinary sentence of $\ipra$.

If the presentation certifies simplicity, the stage areas and their limit are
the geometric surface area by \cref{thm:polygon-pappus}.  Without that
certificate the same proof yields area with multiplicity, which is often the
more stable constructive notion.
\end{proof}

\subsection{Lamina presentations}

For a finite cell union $E$, write $A(E),M(E),V(E)$ for the quantities of
\cref{sec:finite-volume}.  For the nested inner and outer approximants put
\begin{align*}
 A_n^-&=A(I_n),& A_n^+&=A(O_n),\\
 M_n^-&=M(I_n),& M_n^+&=M(O_n),\\
 V_n^-&=V(I_n),& V_n^+&=V(O_n).
\end{align*}
Monotonicity of finite sums and $0\leq r\leq R$ give
\begin{align}
 0&\leq A_n^+-A_n^-\leq2^{-n},
 \label{eq:area-bracket-gap}\\
 0&\leq M_n^+-M_n^-\leq R2^{-n},
 \label{eq:moment-bracket-gap}\\
 0&\leq V_n^+-V_n^-
   =\tau(M_n^+-M_n^-)
   \leq8R2^{-n}.
 \label{eq:volume-bracket-gap}
\end{align}
The equality in \eqref{eq:volume-bracket-gap} is the finite grid theorem
applied to the cell difference after refining both unions to a common grid.

Choose a natural number $B_R$ with $R\leq2^{B_R}$; it can be supplied with the
bounding box or computed once from its rational code.  Then
$n=k+B_R+4$ makes the moment and volume gaps at most $2^{-k}$.  Therefore the
three pairs of brackets define regular names $A,M,V$.

\begin{proof}[Proof of \cref{thm:volume-main}]
For each stage, \cref{thm:grid-pappus} gives both
\[
 V_n^-=\tau M_n^-,\qquad V_n^+=\tau M_n^+.
\]
The nested brackets and \cref{eq:moment-bracket-gap,eq:volume-bracket-gap}
therefore squeeze their common limits to
\[
 V=\tau M.
\]
The supplied positive lower witness for $A$ defines
$\bar r=M/A$ and proves $A\bar r=M$.  Hence
\[
 V=\tau M=\tau A\bar r.
\]
Putting $d=\tau\bar r$ gives $V=Ad$.

For a requested precision $k$, the object proof uses the single explicit
stage $n=k+B_R+7$, refines $I_n,O_n$ to a common dyadic grid, checks the cell
identities by polynomial normalization, and applies
\eqref{eq:area-bracket-gap}--\eqref{eq:volume-bracket-gap}.  There is no
quantification over subsets of the plane and no use of a measure-completion
principle.
\end{proof}

\section{Recovery of the familiar analytic theorem}
\label{sec:bridges}

The main theorems deliberately take effective approximation data as their
hypotheses.  This section explains in detail how standard smooth curves and
Jordan laminas provide those data.  These are bridge theorems: the finite
identities above are the proof-theoretic core, while the bridges identify the
represented limits with customary analytic quantities.

\subsection{Piecewise \texorpdfstring{$C^1$}{C1} curves with explicit moduli}

Consider one parameter interval $[0,1]$ and a curve
\[
 \gamma(t)=(r(t),z(t)),\qquad0\leq r(t)\leq R.
\]
For a direct arithmetization, assume the presentation provides:
\begin{enumerate}[label=(\alph*),leftmargin=2.3em]
\item regular names for $\gamma$ and a derivative candidate
$D\gamma(t)$ at dyadic $t$;
\item rational bounds $\norm{D\gamma(t)}\leq B$ and $0\leq r(t)\leq R$;
\item a primitive-recursive function $\eta(n)$ tending to $0$ with a known
modulus, such that for dyadic $t,h\geq0$ with $t+h\leq1$ and $h\leq2^{-n}$,
\begin{equation}
 \norm{\gamma(t+h)-\gamma(t)-hD\gamma(t)}\leq h\eta(n);
 \label{eq:taylor-modulus}
\end{equation}
\item moduli $\omega_v,\omega_r,\omega_{rv}$ for the bounded functions
$v(t)=\norm{D\gamma(t)}$, $r(t)$, and $r(t)v(t)$.
\end{enumerate}
A usual effectively $C^1$ curve supplies these data: uniform continuity of the
derivative gives \eqref{eq:taylor-modulus}, with the integral remainder proved
by finite Riemann sums.  Stating the remainder estimate directly is more
faithful to first-order arithmetic than treating ``$C^1$'' as an unanalyzed
higher-type predicate.

Take the uniform mesh $t_i=i2^{-n}$ and first use the exact points
$\gamma(t_i)$.  Let $h=2^{-n}$, let $\ell_i$ be the $i$th chord length, and
put $v_i=v(t_i)$.  The reverse triangle inequality applied to
\eqref{eq:taylor-modulus} yields
\begin{equation}
 \abs{\ell_i-hv_i}\leq h\eta(n).
 \label{eq:chord-speed}
\end{equation}
Summing over $2^n$ intervals gives
\begin{equation}
 \left|\sum_i\ell_i-h\sum_i v_i\right|\leq\eta(n).
 \label{eq:length-rsum}
\end{equation}
The Riemann-sum modulus for $v$ then gives
\begin{equation}
 \abs{L(P_n)-L}\leq\eta(n)+\omega_v(n),
 \qquad L=\int_0^1v(t)\,dt,
 \label{eq:length-convergence}
\end{equation}
where the integral notation abbreviates the regular real defined by these
finite sums.

For the moment, write
\[
 m_i=\frac{r(t_i)+r(t_{i+1})}{2}\ell_i.
\]
Using $r\leq R$, $v\leq B$, \eqref{eq:chord-speed}, and
$\abs{r(t_{i+1})-r(t_i)}\leq\omega_r(n)$, one obtains
\begin{align}
 \abs{m_i-hr(t_i)v_i}
 &\leq R\abs{\ell_i-hv_i}
   +hv_i\frac{\abs{r(t_{i+1})-r(t_i)}}2
   +\eta(n)h\frac{\omega_r(n)}2\notag\\
 &\leq h\left(R\eta(n)+\frac{B\omega_r(n)}2
   +\frac{\eta(n)\omega_r(n)}2\right).
 \label{eq:moment-cell-error}
\end{align}
After summation and comparison with the Riemann sum for $rv$,
\begin{equation}
 \abs{M(P_n)-M}
 \leq R\eta(n)+\frac{B\omega_r(n)}2
      +\frac{\eta(n)\omega_r(n)}2+\omega_{rv}(n),
 \qquad M=\int_0^1r(t)v(t)\,dt.
 \label{eq:moment-convergence}
\end{equation}
Every term on the right has a supplied primitive-recursive convergence
modulus.  Hence the polygonal presentation satisfies the surface theorem's
hypotheses.

The vertices used internally must have rational coordinates.  At mesh level
$n$, approximate each coordinate of $\gamma(t_i)$ within
$\delta_n=2^{-(3n+C)}$, where the fixed $C$ dominates $R+B+1$.  Perturbing one
chord length changes it by at most $2\sqrt2\delta_n$; over $2^n$ chords the
total length error is at most $2^{n+2}\delta_n$.  The moment error has the
same form with a fixed factor depending on $R$ and $B$.  These errors are
bounded by $2^{-n}$ after increasing $C$.  Thus rationalization preserves the
explicit Cauchy moduli.

For a finite piecewise-$C^1$ curve, apply the construction on each patch and
concatenate the finite chains.  A fixed finite maximum of the patch moduli is
primitive recursive.  Equations \eqref{eq:length-convergence} and
\eqref{eq:moment-convergence}, followed by the surface theorem, give
\begin{equation}
 S=\tau\int_0^1r(t)\norm{\gamma'(t)}\,dt
  =\tau L\bar r.
 \label{eq:analytic-surface}
\end{equation}
Under the usual injectivity hypothesis this parametric area is the geometric
surface area.  Constructively, that last identification is proved by the same
inner/outer polyhedral squeeze, not by invoking a general area formula as an
axiom.

\subsection{Jordan laminas with a boundary-cover modulus}

Let $D$ lie in a rational box $[0,R]\times[-H,H]$.  The effective datum needed
by \cref{thm:volume-main} is the following.

\begin{definition}[Boundary-cover presentation]
\label{def:boundary-cover}
A primitive-recursive boundary-cover presentation of $D$ consists, for each
$n$, of:
\begin{enumerate}[label=(\roman*),leftmargin=2.3em]
\item a dyadic mesh size and a finite list $B_n$ of mesh squares covering the
boundary, with total area at most $2^{-n}$;
\item for every mesh square not in $B_n$, a decidable label
\code{inside} or \code{outside};
\item coherence certificates saying that an inside square is contained in
$D$, an outside square is disjoint from the interior of $D$, and the labels
are preserved by refinement away from later boundary covers.
\end{enumerate}
\end{definition}

Let $I_n$ be the union of inside squares and let $O_n$ be $I_n$ together with
the boundary-cover squares.  After replacing stage $n$ by the common
refinement of stages $0,\ldots,n$, which is a bounded finite operation, one
gets
\[
 I_n\subseteq I_{n+1}\subseteq O_{n+1}\subseteq O_n,
 \qquad A(O_n\setminus I_n)\leq2^{-n}.
\]
Thus \cref{def:boundary-cover} produces exactly the presentation used in the
volume theorem.

A piecewise-$C^1$ Jordan boundary with effective data supplies such covers.
For example, if a parametrization has speed at most $B$, choose a parameter
mesh fine enough that each image piece has diameter at most $\delta$.  Each
piece meets at most a fixed number of dyadic squares of side comparable to
$\delta$, while the number of pieces is at most a fixed multiple of
$B/\delta+1$.  The total cover area is therefore bounded by
$C(B+1)\delta$.  Choosing
\[
 \delta\leq\frac{2^{-n}}{C(B+1)}
\]
gives the required area bound.  An effective separation or winding-number
certificate labels every square outside the cover.  For polygonal Jordan
curves the label is the parity of a rational ray-crossing count, with all
boundary cases already removed by $B_n$.  For smooth effective Jordan curves,
a polygonal approximation and a supplied separation modulus transfer the same
label.

The inner/outer moment and volume gaps are then exactly
\cref{eq:moment-bracket-gap,eq:volume-bracket-gap}.  In ordinary notation,
\[
 A=\iint_D1\,dr\,dz,
 \qquad
 M=\iint_Dr\,dr\,dz,
 \qquad
 V=\tau M,
\]
but the double integrals here denote the common limits of explicit cell sums.
No Fubini theorem or general change-of-variables theorem is needed for the
proof.

\subsection{Choice of coordinates}

The axis is not selected by a nonconstructive existence argument.  It is part
of the input: an origin point on the axis and an explicit orthonormal frame are
given.  Rational or regular-real coordinate transformations then reduce the
problem to the meridional coordinates $(r,z)$.  The radial coordinate is
nonnegative by a supplied half-plane certificate.  Consequently all later
uses of ``distance from the axis'' are fixed primitive-recursive operations on
presented coordinates.

\section{Formal status inside intuitionistic primitive recursive arithmetic}
\label{sec:ipra-audit}

This section records why the preceding proof is genuinely available in
$\ipra$, rather than merely constructive in an informal sense.

\subsection{The first-order theorem schema}

Fix one surface or lamina presentation.  Its algorithms are named by function
symbols introduced through primitive-recursive definitions.  The assumptions
about those functions are universal closures of equations and decidable
inequalities.  For example, real equality is represented by
\begin{equation}
 X\Eq Y
 \quad:\Longleftrightarrow\quad
 \forall k\;
 \abs{m_X(k+3)-m_Y(k+3)}\leq2^{-k},
 \label{eq:real-equality-code}
\end{equation}
where $m_X(k)$ is a rational midpoint code and the displayed inequality is
implemented by cross multiplication.  Equivalent overlap formulations using
$L_X,U_X,L_Y,U_Y$ are interchangeable with a fixed precision shift.

Thus the final surface sentence has the form
\[
 \forall k\;E_{\mathcal C}(k),
\]
where $E_{\mathcal C}$ is a quantifier-free formula asserting the precision
$k$ instance of
$S=\tau M=\tau L(M/L)$.  The volume sentence is analogous.  There is one such
sentence for each fixed primitive-recursive presentation.  This is exactly how
ordinary $\ipra$ handles any theorem about an algorithm: it quantifies over
natural-number inputs, not over a completed function space.

\subsection{Inventory of primitive-recursive constructions}

The proof uses only the following closure properties.
\begin{longtable}{@{}p{0.28\textwidth}p{0.66\textwidth}@{}}
\toprule
Construction & Primitive-recursive realization \\
\midrule
Rationals and dyadics & Signed numerator, positive denominator, and exponent
codes; all order tests by cross multiplication.\\
Finite sequences & Pairing code, bounded length, projection, map, fold, sum,
product, maximum, and common-grid refinement.\\
Square root & Bounded bisection in a supplied rational interval; the number of
steps is the requested precision plus a fixed magnitude bound.\\
Circle constants & The recurrence for $c_m$, bounded square roots, and finite
field operations; all error estimates are explicit in $m$.\\
Segment lengths & Square root of the rational sum of two squares, represented
by bisection intervals.\\
Polyhedral areas & Finite sums of trapezoid areas from
\cref{eq:inner-frustum,eq:outer-frustum}.\\
Grid areas and moments & Bounded folds of rational polynomial expressions.\\
Limits & Diagonal selection at a stage computed from the supplied Cauchy
modulus.\\
Reciprocal & Interval reversal after a supplied positive lower witness and a
fixed precision shift.\\
Boundary approximants & Bounded decoding, sorting, duplicate removal, cell
classification, and finite union refinement.\\
\bottomrule
\end{longtable}

Each item is built from composition and primitive recursion.  Any bounded
search appearing in bisection or cell processing has a numerical bound
computed before the search begins.

\subsection{Induction ledger}

Only quantifier-free induction is required.
\begin{enumerate}[leftmargin=2.2em]
\item Induction on a list length proves correctness of finite sums, moments,
and the additivity identities.
\item Induction on the bisection count proves the square-root interval
invariant.
\item Induction on $m$ proves the circle recurrences and the estimates
\eqref{eq:c-bound}--\eqref{eq:area-gap}.
\item Induction on the number of polygon segments proves
\cref{thm:polygon-pappus}; induction on the number of cells proves
\cref{thm:grid-pappus}.
\item Induction on refinement depth proves nesting of the inner and outer
lamina codes.
\end{enumerate}
All predicates in these inductions are Boolean combinations of equalities and
orders of coded integers.  Existential witnesses such as a bisection midpoint,
a precision stage, or a common refinement are given by primitive-recursive
terms before induction begins.

\subsection{Principles not used}

The proof does not use:
\begin{itemize}[leftmargin=2em]
\item excluded middle for arbitrary formulas or double-negation elimination;
\item countable, dependent, or unique choice;
\item an axiom asserting that every Cauchy sequence has a real limit;
\item arbitrary suprema, infima, measurable-set completion, or sigma-additivity;
\item a general Fubini, Tonelli, coarea, or change-of-variables theorem;
\item a nonconstructive theorem choosing an arc-length parametrization;
\item a decision procedure for equality of arbitrary represented reals.
\end{itemize}
Every limit used is a particular diagonal sequence whose modulus is already
known.  Every division is accompanied by a positive lower witness.

\begin{remark}[Why positivity is explicit]
Classically, a nonzero compact curve has positive length and a nonempty Jordan
region has positive area under familiar hypotheses.  Constructively, a proof
of $X\neq0$ does not in general compute a rational $\lambda>0$ with
$X\geq\lambda$.  The latter is what reciprocal formation needs.  Requiring the
witness is therefore substantive and mathematically honest.  Standard
polygonal, smooth, and Jordan inputs normally provide it directly.
\end{remark}

\section{Lean formalization plan}
\label{sec:lean}

The repository snapshot used for this plan is
\[
 \code{VladimirReshetnikov/ProveIt@8c135f6f2a9d041cde785a1a632c44ad98065b20}.
\]
The existing development already supplies two complementary layers.

\begin{itemize}[leftmargin=2em]
\item \code{ConciseIPRA.PR} defines primitive-recursive symbols
\code{PR}, ordinary $\ipra$ terms \code{PTm}, the axiom theory
\code{PRAAx}, its recursion equations, and the theorem \code{PRA.qfInd}.
\item \code{ConciseIPRA.PRAToScalar} compiles every \code{PR} symbol to the
native iterator term language and packages the result as
\code{FromPRA.interp : Interp PRAAx TeqAx}.
\item \code{ConciseIPRA.Interpretation} proves \code{Interp.sound}, so an
ordinary $\ipra$ derivation transports to the concise scalar theory.
\item \code{ConciseIPRA.BiInterpretation} exposes
\code{pra_iff_teq} and \code{pra_theorem_iff}; these can be used as the final
faithfulness bridge.
\end{itemize}

The recommended strategy separates semantic mathematics from object-theory
arithmetization.  This makes every intermediate theorem testable in ordinary
Lean while keeping the final claim genuinely about derivability in $\ipra$.

\subsection{Phase I: executable semantic library}

Create the following modules under a new \code{Pappus} application directory.

\begin{longtable}{@{}p{0.31\textwidth}p{0.63\textwidth}@{}}
\toprule
Module & Main content and terminal lemmas \\
\midrule
\code{Code.lean} & Signed rationals, dyadics, bounded lists, cell and vertex
codes, decoding correctness.\\
\code{IntervalReal.lean} & Regular interval names, equality, bounded
arithmetic, diagonal limits, and positive reciprocal.\\
\code{Sqrt.lean} & Rational bisection and proofs of nesting, width, and
$s^2=x$ at represented-real precision.\\
\code{Circle.lean} & $c_m,s_m,p_m,q_m,a_m,b_m$ and the estimates of
\cref{lem:polygon-estimates}; definitions of \code{tau} and \code{pi}.\\
\code{Segment.lean} & Segment length and midpoint moment; inner and outer
polyhedral frustum brackets; \code{segmentSweep_eq_tau_mul_moment}.\\
\code{Polygon.lean} & Finite folds for length, moment, and sweep;
\code{polygonSweep_eq_tau_mul_moment}; simplicity certificate interface.\\
\code{Cell.lean} & Rectangle area, moment, annular-cylinder volume, and
\code{cellVolume_eq_tau_mul_moment}.\\
\code{GridLamina.lean} & Common-grid refinement, disjointness, finite
additivity, and \code{gridVolume_eq_tau_mul_moment}.\\
\code{LimitTransfer.lean} & The quantitative identity-transfer lemma and
precision-shift combinators.\\
\code{Surface.lean} & Presented polygonal sequences and the semantic theorem
\code{surface_pappus}.\\
\code{Volume.lean} & Nested inner/outer cell presentations and the semantic
theorem \code{volume_pappus}.\\
\code{SmoothBridge.lean} & Chord and moment estimates
\eqref{eq:chord-speed}--\eqref{eq:moment-convergence}.\\
\code{JordanBridge.lean} & Boundary-cover presentations and construction of
nested cell unions.\\
\bottomrule
\end{longtable}

The finite algebra file included in this archive, \path{PappusFinite.lean}, is
the first slice of \path{Segment.lean}, \path{Polygon.lean}, \path{Cell.lean},
and \path{GridLamina.lean}.  It proves the list-level Pappus identities over
an arbitrary linear ordered field.  It intentionally treats the segment sweep
formula as the semantic result supplied by the later frustum-bracket module;
the nontrivial next step is to connect that formula to the inner and outer
polyhedra.

A representative semantic theorem will have a shape close to:
\begin{lstlisting}
theorem polygon_pappus
    (P : List (SegmentDatum K)) :
    polySweep tau P = tau * polyMoment P := by
  induction P with
  | nil => simp [polySweep, polyMoment]
  | cons s P ih =>
      simp [polySweep, polyMoment, SegmentDatum.sweep, ih, mul_add]
\end{lstlisting}
The actual represented-real version replaces field equality by
\code{RealName.Eq} and invokes the segment bracket theorem.

\subsection{Phase II: primitive-recursive definitions}

For every executable operation used by the fixed presentation, define a symbol
of the inductive family \code{PR k}.  Composition and primitive recursion are
already constructors, so this is mostly a library of named abbreviations.
Suggested modules are:
\begin{itemize}[leftmargin=2em]
\item \code{PRCode.lean}: pairing, signed integers, rationals, dyadics, lists;
\item \code{PRBisection.lean}: bounded square-root bisection;
\item \code{PRGeometry.lean}: segment, cell, polygon, and grid folds;
\item \code{PRPresentations.lean}: the particular surface and lamina
presentation symbols and their moduli.
\end{itemize}
For example, a list fold is a \code{PR.prec} whose base is zero and whose step
adds the decoded next contribution.  The equations needed later are direct
instances of \code{PRA.precZero} and \code{PRA.precSucc}.

\subsection{Phase III: object-level proof graphs}

Define quantifier-free graph formulas for the coded operations.  Prove, in
\code{PRAAx}, the following families.
\begin{enumerate}[leftmargin=2.2em]
\item Functionality and defining equations of rational arithmetic and finite
folds.
\item Bisection invariants and the square-root width theorem.
\item Circle polygon estimates and the bound $\tau<8$.
\item Denominator-cleared segment and cell identities.
\item Cauchy and nesting invariants supplied by each fixed presentation.
\item Precision-level versions of the two limit-transfer arguments.
\end{enumerate}
Use \code{PRA.qfInd} whenever an index is inductive.  Since all graph formulas
are quantifier-free after witnesses have been named by terms, no stronger
induction is required.

A schematic final object theorem is:
\begin{lstlisting}
def SurfacePappusQF (k : PTm n) : Fm PTm n :=
  -- denominator-cleared precision-k overlap of
  -- S, tau * M, and tau * L * (M / L)

 theorem surface_pappus_pra :
   [] |-[PRAAx] Fm.all SurfacePappusQF := by
   apply PRA.qfInd SurfacePappusQF (by decide)
   case zero => exact surface_base
   case succ => exact surface_step
\end{lstlisting}
The archive roadmap replaces the ASCII turnstile by the actual notation
imported from \path{ConciseIPRA.PR}; the display above emphasizes the proof
shape rather than committing this report to a parser-sensitive glyph.

\subsection{Phase IV: transport to the concise scalar theory}

Once \code{surface_pappus_pra} and \code{volume_pappus_pra} are complete,
transport them with
\begin{lstlisting}
def surface_pappus_teq :=
  FromPRA.interp.sound surface_pappus_pra

def volume_pappus_teq :=
  FromPRA.interp.sound volume_pappus_pra
\end{lstlisting}
or obtain the equivalent statements using \code{pra_theorem_iff}.  This step
is small because the repository already proves that the interpretations
preserve derivability and that the two theories are bi-interpretable.

\subsection{Completion criteria}

A complete Lean formalization should satisfy all of the following.
\begin{enumerate}[leftmargin=2.2em]
\item The semantic modules compile without \code{sorry}, classical choice, or
noncomputable definitions in the theorem dependency cone.
\item \code{#print axioms} for the semantic Pappus theorems shows only Lean's
constructive kernel and the explicitly chosen quotient infrastructure, if any.
A quotient-free interval-name implementation is preferable.
\item The object-level theorems are terms of
\code{[] |-[PRAAx] ...}; they are not merely semantic theorems about
\code{Nat}.
\item Translation to \code{TeqAx} is checked through the existing
interpretation API.
\item Executable tests compare the object-level functions with the independent
Python checker on many small polygon and cell certificates.
\end{enumerate}

\section{Worked finite examples}
\label{sec:examples}

\subsection{A polygonal generator}

Take the chain
\[
 (1,0)\longrightarrow(4,4)\longrightarrow(4,7).
\]
The segment lengths are $5$ and $3$.  Hence
\[
 L=8,
 \qquad
 M=5\frac{1+4}{2}+3\frac{4+4}{2}
   =\frac{49}{2},
 \qquad
 \bar r=\frac{49}{16}.
\]
The swept surface area is therefore
\[
 S=\frac{49}{2}\tau=8\tau\frac{49}{16}.
\]
The accompanying JSON certificate contains these data.  The Python checker
verifies the two Pythagorean length equations and every rational identity
exactly, using \code{fractions.Fraction}.

\subsection{An L-shaped grid lamina}

Take the two cells
\[
 [1,3]\times[0,2],
 \qquad
 [3,4]\times[0,1].
\]
They have disjoint interiors.  Their total area and moment are
\[
 A=4+1=5,
 \qquad
 M=2\cdot4+\frac72\cdot1=\frac{23}{2},
 \qquad
 \bar r=\frac{23}{10}.
\]
The rotated volume is
\[
 V=\frac{23}{2}\tau=5\tau\frac{23}{10}.
\]
Again the checker verifies both the shell formula and the centroid identity by
integer cross multiplication.

\subsection{The torus as a limit example}

Let a circle of radius $a$ have center at radial distance $R>a$ from the axis.
Its polygonal circle presentations satisfy the hypotheses of both main
theorems.  The generating circumference is $\tau a$ and its centroid has
radius $R$, so the surface formula gives
\[
 S=\tau^2Ra=4\pi^2Ra.
\]
The disk has area $\pi a^2$ and the same centroid radius, so the volume formula
gives
\[
 V=\tau\pi Ra^2=2\pi^2Ra^2.
\]
These equalities are corollaries of the same regular-polygon circle package;
no independent transcendental theory of $\pi$ is required.

\section{Scope, variants, and edge cases}
\label{sec:scope}

\begin{enumerate}[leftmargin=2.2em]
\item \textbf{Axis contact.}  The classical statement often says that the
axis is external.  The proofs above need only $r\geq0$ and permit contact on
boundary pieces of zero relevant content.  A generator segment lying along
the axis contributes zero surface area.  A lamina cell of zero radial width
contributes zero volume.

\item \textbf{Self-overlap.}  For arbitrary polygonal chains the surface
identity is exact with multiplicity.  For an ordinary geometric-area statement,
a simplicity or no-retracing certificate is necessary.  The volume theorem's
nested cell unions are sets by construction and therefore need no multiplicity
convention.

\item \textbf{Density.}  A weighted version follows immediately.  Replace
length or area by a nonnegative primitive-recursive density and define the
corresponding weighted radial moment.  Rotating each infinitesimal mass does
not by itself yield unweighted surface area or volume, but it yields the
Pappus identity for the swept weighted content.  For piecewise-constant
densities the finite proof is literal; general presented densities follow by
explicit limits.

\item \textbf{Signed and oriented forms.}  If one allows signed radial
coordinates and signed cell multiplicities, the algebraic identities remain
valid without the half-plane restriction.  They then compute oriented surface
or volume moments and automatically cancel overlaps.  Recovering unsigned
geometric content requires positivity and nonoverlap information.

\item \textbf{Zero length or area.}  The identities $S=\tau M$ and
$V=\tau M$ remain valid.  Only the centroid quotients $M/L$ and $M/A$ are
undefined without positive lower witnesses.
\end{enumerate}

\section{Further research questions}
\label{sec:further}

\begin{enumerate}[leftmargin=2.2em]
\item Can the effective Jordan-curve bridge be fully internalized in the
current \code{ConciseIPRA} library, including a compact primitive-recursive
winding-parity classifier and a quantitative separation theorem?

\item What is the smallest object-language circle package sufficient for the
proof?  The regular-polygon construction is transparent, but a Machin-type
rational series might produce shorter certificates even though it obscures
the geometry.

\item Can one prove, in $\ipra$, equivalence between the polygonal-squeeze
surface content used here and a suitably coded two-dimensional Hausdorff
content for every effective simple surface of revolution?

\item How much of the semantic-to-object translation can be automated?  A
verified compiler from equations proved over executable natural-number
functions to \code{PRAAx} derivations would make applications of
\code{PRA.qfInd} substantially shorter.

\item Is there a useful centroid-free formulation that also covers vanishing
length or area uniformly, taking $M$ rather than $\bar r$ as the primitive
quantity?  The proof suggests that $S=\tau M$ and $V=\tau M$ are the logically
stronger and constructively cleaner statements.

\item Can the no-retracing certificate for surfaces be weakened to an
explicit overlap multiplicity bound, producing a constructive area formula
with a computable multiplicity correction?
\end{enumerate}

\section{Conclusion}

Pappus's theorem is often presented as a one-line consequence of surface and
volume integrals.  In the restricted constructive setting, the direction of
explanation reverses.  The primitive facts are finite: a regular polygon, a
planar trapezoid, an annular prism, a midpoint sum, and a bounded list fold.
Those facts prove
\[
 \Surf(\text{segment})=\tau\Mom(\text{segment}),
 \qquad
 \Vol(\text{cell})=\tau\Mom(\text{cell}).
\]
Finite additivity gives the polygonal and grid theorems.  Explicit Cauchy and
boundary-cover moduli then pass the identities to regular real limits.  The
centroid formulas are obtained only at the final step, by dividing the first
moment by a positively witnessed length or area.

This organization exposes exactly what $\ipra$ must prove and exactly where
effective data are required.  It also yields a direct Lean development path:
complete the regular-real and polygon-squeeze semantic modules, arithmetize
their primitive-recursive algorithms, prove the precision-level formulas by
quantifier-free induction, and invoke the existing
\code{ConciseIPRA} interpretation and bi-interpretation machinery.

\appendix

\section{A precision ledger}
\label{app:precision}

The following conservative shifts suffice; sharper constants are unnecessary.
Assume all fixed magnitude constants have been bounded by powers of two in a
preprocessing certificate.

\begin{longtable}{@{}p{0.34\textwidth}p{0.58\textwidth}@{}}
\toprule
Target at precision $2^{-k}$ & Sufficient input precision/stage \\
\midrule
Sum of at most $2^e$ terms & Each term at precision $k+e+2$.\\
Product $XY$, $\abs X,\abs Y\leq2^b$ & Inputs at precision $k+b+3$.\\
Multiplication by $\tau$ & Input at precision $k+4$, using $\tau<8$.\\
Square root on $[0,2^{2b}]$ & Bisection interval width at most
$2^{-(k+b+3)}$; compare squares to maintain nesting.\\
Segment frustum brackets & Choose $m\geq k+B_\sigma+3$, where
$2^{B_\sigma}$ bounds
$4(r_0+r_1)(\ell+\abs{\Delta r})$.\\
Cell-union moment brackets & Choose $n\geq k+B_R+1$ with $R\leq2^{B_R}$.\\
Cell-union volume brackets & Choose $n\geq k+B_R+4$.\\
Positive reciprocal $1/X$, $X\geq2^{-b}$ & Read $X$ at precision
$k+2b+3$.\\
Diagonal limit of $x_n$ & Evaluate $x_{\mu(k+3)}$ at precision $k+3$.\\
Rationalized mesh with $2^n$ vertices & Coordinate error
$2^{-(3n+C)}$, where $C$ bounds the fixed geometry constants.\\
\bottomrule
\end{longtable}

Every row is a primitive-recursive expression in $k$ and the fixed
certificate parameters.  These shifts can therefore be inserted literally in
the object-level formulas, avoiding all metatheoretic phrases such as
``take $n$ sufficiently large.''

\section{Archive contents and reproducibility}
\label{app:archive}

The accompanying archive contains:
\begin{itemize}[leftmargin=2em]
\item \code{pappus_centroid_ipra.tex} and the compiled
\code{pappus_centroid_ipra.pdf};
\item \code{lean/PappusFinite.lean}, a no-\code{sorry} finite algebra scaffold
(not locally kernel-checked);
\item \code{lean/PappusRoadmap.lean}, parser-oriented declarations and theorem
stubs marking the remaining semantic and object-theory obligations;
\item \code{FORMALIZATION_PLAN.md}, an obligation-by-obligation Lean roadmap;
\item \code{code/pappus_checker.py} and
\code{certificates/sample_certificate.json};
\item \code{verification_output.txt}, the checker's recorded output;
\item \code{SOURCE_PROVENANCE.md}, \code{MANIFEST.txt}, and
\code{SHA256SUMS};
\item a \code{Makefile} and \code{build.sh} for rebuilding the PDF, running
the checker, and recreating the checksums.
\end{itemize}

The Python checker is not a substitute for the planned Lean proof.  It is an
independent executable audit of the exact rational finite identities and the
sample certificates.  The Lean scaffold likewise does not claim that the
full represented-real or $\ipra$ derivation is already complete; its purpose
is to make the next formal steps explicit and type-directed.

\begingroup
\small
\begin{thebibliography}{9}

\bibitem{MathWorldPappus}
Eric W. Weisstein,
``Pappus's Centroid Theorem,'' \emph{MathWorld---A Wolfram Web Resource},
\href{https://mathworld.wolfram.com/PappussCentroidTheorem.html}{MathWorld entry}.

\bibitem{ProveItConciseIPRA}
Vladimir Reshetnikov,
\emph{ProveIt: Concise presentations of intuitionistic primitive recursive
arithmetic}, commit
\code{8c135f6f2a9d041cde785a1a632c44ad98065b20}, 15 November 2026,
\href{https://github.com/VladimirReshetnikov/ProveIt/tree/main/Logic/PrimitiveRecursiveArithmetic/ConciseIPRA}{ProveIt ConciseIPRA directory}.

\end{thebibliography}
\endgroup

\end{document}
